\PassOptionsToPackage{dvipsnames}{color}
\documentclass[10pt]{NSP1}
\usepackage{url,floatflt}
\usepackage{helvet,times}
\usepackage{psfig,graphics}
\usepackage{mathptmx,amsmath,amssymb,bm}
\usepackage{float}
\usepackage[bf,hypcap]{caption}

\addtolength{\oddsidemargin}{.5in}
\addtolength{\evensidemargin}{-.5in}

\tolerance=1
\emergencystretch=\maxdimen
\hyphenpenalty=10000
\hbadness=10000

\topmargin=0.00cm

\def\sm{\smallskip}
\def\no{\noindent}

\def\firstpage{1}
\setcounter{page}{\firstpage}
\def\thevol{1}
\def\thenumber{2}
\def\theyear{2015}

\numberwithin{equation}{section}



\begin{document}

\titlefigurecaption{{\large \bf \rm Progress in Fractional Differentiation
and Applications}\\ {\it\small An International Journal}}

\title{49 Amplified Generalized Fractional Neural Network Approximations at
a reduced finite domain}

\author{George A. Anastassiou{$^{1}$}}
\institute{Department of Mathematical Sciences\\
University of Memphis\\
Memphis, TN 38152, U.S.A.}

\titlerunning{49 Amplified Generalized Fractional Neural Network Approximations ...}
\authorrunning{G. A. Anastassiou}

\mail{ganastss@memphis.edu}

\received{}
\revised{}
\accepted{}
\published{}

\abstracttext{This research deals with the determination of the rate of generalized fractional pointwise, uniform and $L_{p}$, $p\geq 1,$ convergences to the unit operator of 49 specific "normalized cusp neural network operators". The cusp is a compact support activation function, which derives from the composition of two specific activation functions having as domain the whole real line. We employ 7 different known activation functions. These convergences are given via the moduli of continuity of the right and left generalized fractional derivatives of Caputo type of the engaged function in the form of Jackson type inequalities. The composition of activation functions aims to more flexible and powerful neural networks utilizing the reduction of infinite domains to the one domain of compact support.}

\keywords{Amplified Neural network approximation, cusp
activation function, modulus of continuity, reduction of domain, generalized
fractional derivatives.\\[2mm]
\textbf{2020 AMS Mathematics Subject Classification:} 26A33, 41A17, 41A25,
41A30, 41A36.}

\maketitle

\section{\; Introduction}

From AI and computer science we have the following: In essence, composing
activation functions in neural networks offers the advantage of potentially
tailoring the network's ability to learn and model complex, non-linear
relationships in data. Here's a breakdown of the potential benefits:

1. Enhanced Capacity for Complex Modeling:

\begin{itemize}
\item Diversification of Non-linearity: Different activation functions have
different characteristics. For example, ReLU introduces sparsity, while
Sigmoid squashes values into a range. By composing them, the network
potentially can learn a wider variety of non-linear transformations and
capture more intricate patterns in the data.
\end{itemize}

2. Improved Training Dynamics:

\begin{itemize}
\item Mitigating Gradient Problems: Activation functions influence gradient
flow during training. Using different activation functions can potentially
help address issues like vanishing or exploding gradients, which hinder
learning in deep networks.

\item Faster Convergence: Certain activation functions, like ReLU, can
accelerate the convergence of the training process compared to others like
Sigmoid or Tanh. Combining different functions can potentially lead to
faster training and competitive performance.
\end{itemize}

3. Enhanced Generalization and Robustness:

\begin{itemize}
\item Better Generalization: By learning richer representations of the data
through diverse activation functions, the network's ability to generalize
well to unseen data improves, reducing the risk of overfitting.

\item Increased Robustness: Networks with carefully chosen activation
functions can handle variations in input data more effectively, adapting to
noise, missing data, or unexpected perturbations.
\end{itemize}

4. Adaptation to Input Characteristics:

\begin{itemize}
\item Handling Diverse Data: Different activation functions can be suited to
different data characteristics. For instance, tanh can be useful when
dealing with data containing both positive and negative values.
\end{itemize}

5. Potential for Architectural Interpretability:

\begin{itemize}
\item Insight into Learning: By using distinct activation functions,
different parts of the network might become responsible for capturing
specific features, which can potentially offer insights into how the model
learns.
\end{itemize}

In summary, composing activation functions potentially allows for a more
flexible and powerful neural network capable of:

\begin{itemize}
\item Learning more complex patterns.

\item Faster and more stable training.

\item Better generalization to new data.

\item Greater adaptability to diverse data.
\end{itemize}

Attention: While composing activation functions can offer benefits, it's
important to choose them judiciously and with consideration for the specific
problem at hand, as some combinations might not be beneficial or could even
lead to unwanted behaviors like exploding gradients. Empirical testing and
validation are crucial when exploring different activation function
compositions.

The author was greatly inspired and motivated by \cite{8} and he was the
pioneer of quantitative neural network approximation, see \cite{2}, and
since then he has published numerous of papers and books, e.g. see \cite{3}-
\cite{7}.

In this article we continue this trend at the generalized fractional level.

In mathematical neural network approximation AMS Mathscinet lists no
articles related to composition of activation functions. So this is a
seminal article.

By using composition of known and specific activation functions he achieves
the first long part of the introduction, and most notably these compositions
lead to activation functions of compact support, though the initial
activation functions had an infinite domain, the whole real line.

Now each resulting bi-composed activation function is an open cusp of
compact support $\left[ -1,1\right] $. Our involved 7 activation functions
generate 49 amplified neural network operators, which resemble to the
squashing operators in \cite{2}, \cite{4}, and so do the produced
quantitative results.

Of great inspiration are the articles \cite{12}-\cite{14}. References \cite
{15}, \cite{16}, \cite{18} are foundational. Finally references \cite{19}-
\cite{28} represent recent important works.

\section{\; Foundations}

We need

\begin{definition}
\label{d1}(\cite{7}, p. 530) The sigmoid activation functions we deal with
here are described as follows $g:\mathbb{R}\rightarrow \left[ -1,1\right] ,$
strictly increasing $g\left( 0\right) =0$, $g\left( -x\right) =-g\left(
x\right) $, $x\in \mathbb{R}$, $g\left( +\infty \right) =1$, $g\left(
-\infty \right) =-1$. Also $g$ is strictly convex over $(-\infty ,0]$ and
strictly concave over $[0,+\infty )$, with $g^{\left( 2\right) }\in C\left(
\mathbb{R}\right) $.
\end{definition}

More precisely we will deal with the following sigmoid activation functions
which possess all of the above properties: the algebraic activation
function,
\begin{equation}
g_{1}\left( x\right) =\frac{x}{\sqrt[2m]{1+x^{2m}}}\text{, \ }m\in \mathbb{N}
\text{, }x\in \mathbb{R},  \tag{1}  \label{1.}
\end{equation}
(see \cite{6}, pp. 2-3);

the normalized and parametrized arctangent
\begin{equation}
g_{2}\left( x\right) =\frac{2}{\pi }\arctan \left( \frac{x}{2}\lambda
x\right) =\frac{2}{\pi }\int_{0}^{\frac{\pi \lambda x}{2}}\frac{dz}{1+z^{2}},
\text{ \ }x\in \mathbb{R}\text{, }\lambda >0,  \tag{2}  \label{2.}
\end{equation}
(see \cite{7}, p.p. 156-157);

the normalized and parametrized Gudermannian function
\begin{equation*}
g_{3}\left( x\right) =\frac{2}{\pi }gd\left( \lambda x\right) =\frac{4}{\pi }
\arctan \left( \tanh \left( \frac{\lambda x}{2}\right) \right)
\end{equation*}
\begin{equation}
=\frac{2}{\pi }\int_{0}^{\lambda x}\frac{dt}{\cosh t}=\frac{4}{\pi }
\int_{0}^{\lambda x}\frac{dt}{e^{t}+e^{-t}},\text{ \ }x\in \mathbb{R}
\text{, }\lambda >0,  \tag{3}  \label{3.}
\end{equation}
where
\begin{equation*}
gd\left( x\right) :=\int_{0}^{x}\frac{dt}{\cosh t}=2\arctan \left( \tanh
\left( \frac{x}{2}\right) \right) ,\text{ \ }\forall \text{ }x\in \mathbb{R},
\end{equation*}
is the Gudermannian function (see \cite{7}, p. 198-199);

the parametrized error function
\begin{equation}
g_{4}\left( x\right) =erf\lambda x=\frac{2}{\pi }\int_{0}^{\lambda
x}e^{-t^{2}}dt,\text{ \ }x\in \mathbb{R}\text{, }\lambda >0,  \tag{4}
\label{4.}
\end{equation}
(see \cite{7}, pp. 226-227);

the parametrized hyperbolic tangent function
\begin{equation*}
g_{5}\left( x\right) =\tanh \lambda x=\frac{e^{\lambda x}-e^{-\lambda x}}{
e^{\lambda x}+e^{-\lambda x}}=\frac{e^{2\lambda x}-1}{e^{2\lambda x}+1}
\end{equation*}
\begin{equation}
=\frac{1-e^{-2\lambda x}}{1+e^{-2\lambda x}},\text{ \ }x>0\text{, }\lambda >0
\text{,}  \tag{5}  \label{5.}
\end{equation}
(see \cite{7}, p. 72);

for small $0<\lambda <1$, $\tanh \lambda x$ is expected to behave better
than $ReLu$ and $Leaky$ $ReLu$ activation functions;

the parametrized hyperbolic tangent like activation function,
\begin{equation}
g_{6}\left( x\right) =\frac{A^{\lambda x}-A^{-\lambda x}}{A^{\lambda
x}+A^{-\lambda x}},\text{ \ }A>1\text{, }\lambda >0\text{, }x\in \mathbb{R}
\tag{6}  \label{6.}
\end{equation}
(see \cite{7}, pp. 262-263);

finally, the $\beta $-parametrized half hyperbolic tangent function
\begin{equation}
g_{7}\left( x\right) =\frac{1-e^{-\beta x}}{1+e^{-\beta x}},\text{ \ }\beta
>0\text{, }\forall \text{ }x\in \mathbb{R},  \tag{7}  \label{7.}
\end{equation}
(see \cite{7}, p. 509-510).

So, we can denote then by $g_{i}\left( x\right) $, $x\in \mathbb{R}$, $%
i=1,2,...,7.$

Here we are concerned with the compositions of these activation functions
\begin{equation}
g_{i}\circ g_{j}=g_{i}|_{\left( -1,1\right) }\circ g_{j},\text{ \ and }
g_{i}\circ g_{i}=g_{i}|_{\left( -1,1\right) }\circ g_{i},  \tag{8}
\label{8.}
\end{equation}
where $i,j\in \left \{ 1,2,...,7\right \} .$

There are $P\left( 7,2\right) $ permutations, i.e.
\begin{equation*}
P\left( 7,2\right) =\frac{7!}{5!}=6\cdot 7=42,\text{ of ordered compositions;
}
\end{equation*}
plus $7$ self-compositions; so we have a total of $49$ total such
compositions.

So here a $g_{i}\circ g_{j}$ means any of the compositions in (\ref{8.}),
where $i,j\in \left \{ 1,...,7\right \} $, $i\neq j$ or $i=j.$

Clearly, $g_{i}\circ g_{j}=g_{i}|_{\left( -1,1\right) }\circ g_{j}$ is
strictly increasing and $\left( g_{i}\circ g_{j}\right) \left( 0\right) =0$,
and
\begin{equation*}
g_{i}\circ g_{j}\left( -x\right) =g_{i}\left( g_{j}\left( -x\right) \right)
=g_{i}\left( -g_{j}\left( x\right) \right) =-g_{i}\left( g_{j}\left(
x\right) \right) =-\left( g_{i}\circ g_{j}\right) \left( x\right) ,
\end{equation*}
that is
\begin{equation*}
\left( g_{i}\circ g_{j}\right) \left( -x\right) =-\left( g_{i}\circ
g_{j}\right) \left( x\right) ,\text{ \ }\forall \text{ }x\in \mathbb{R}.
\end{equation*}
Furthermore
\begin{equation*}
\left( g_{i}\circ g_{j}\right) \left( +\infty \right) =g_{i}\left(
g_{k}\left( +\infty \right) \right) =g_{i}\left( 1\right) ,
\end{equation*}
\begin{equation*}
\left( g_{i}\circ g_{j}\right) \left( -\infty \right) =g_{i}\left(
g_{j}\left( -\infty \right) \right) =g_{i}\left( -1\right) .
\end{equation*}
Next acting over $(-\infty ,0]:$ let $\lambda ,\mu \geq 0:\lambda +\mu =1$.
Then by convexity of $g_{j}$ there we have
\begin{equation*}
g_{j}\left( \lambda x+\mu y\right) \leq \lambda g_{j}\left( x\right) +\mu
g_{j}\left( y\right) ,\text{ \ }x,y\in \mathbb{R}_{-};
\end{equation*}
and
\begin{equation*}
g_{i}\left( g_{j}\left( \lambda x+\mu y\right) \right) \leq g_{i}\left(
\lambda g_{j}\left( x\right) +\mu g_{j}\left( y\right) \right) \leq \lambda
g_{i}\left( g_{j}\left( x\right) \right) +\mu g_{i}\left( g_{j}\left(
y\right) \right) ,
\end{equation*}
i.e.
\begin{equation*}
\left( g_{i}\circ g_{j}\right) \left( \lambda x+\mu y\right) \leq \lambda
\left( g_{i}\circ g_{j}\right) \left( x\right) +\mu \left( g_{i}\circ
g_{j}\right) \left( y\right) ,
\end{equation*}
$x,y\in \mathbb{R}_{-}.$

So that $g_{i}\circ g_{j}$ is convex over $(-\infty ,0].$

Similarly, over $[0,+\infty )$ we get: let $\lambda ,\mu \geq 0:\lambda +\mu
=1$. Then by concavity of $g_{j}$ there we have
\begin{equation*}
g_{j}\left( \lambda x+\mu y\right) \geq \lambda g_{j}\left( x\right) +\mu
g_{j}\left( y\right) ,\text{ \ }x,y\in \mathbb{R}_{+};
\end{equation*}
and
\begin{equation*}
g_{i}\left( g_{j}\left( \lambda x+\mu y\right) \right) \geq g_{i}\left(
\lambda g_{j}\left( x\right) +\mu g_{j}\left( y\right) \right) \geq \lambda
g_{i}\left( g_{j}\left( x\right) \right) +\mu g_{i}\left( g_{j}\left(
y\right) \right) .
\end{equation*}
Therefore $g_{i}\circ g_{j}$ is concave over $[0,+\infty ).$

Also, it is
\begin{equation*}
(g_{i}\left( g_{j}\left( x\right) \right) )^{\prime \prime }=g_{i}^{\prime
\prime }\left( g_{j}\left( x\right) \right) \left( g_{j}^{\prime }\left(
x\right) \right) ^{2}+g_{i}^{\prime }\left( g_{j}\left( x\right) \right)
g_{j}^{\prime \prime }\left( x\right) \in C\left( \mathbb{R}\right) \text{,
\ }x\in \mathbb{R}.
\end{equation*}
So $g_{i}\circ g_{j}$ is a sigmoid activation function with asymptotes $%
g_{i}\left( \pm 1\right) $.

Next we consider the function
\begin{equation*}
\Phi _{i,j}\left( x\right) :=\frac{1}{4}\left( g_{i}\circ g_{j}\left(
x+1\right) -g_{i}\circ g_{j}\left( x-1\right) \right) >0,\text{ \ }\forall
\text{ }x\in \mathbb{R}.
\end{equation*}
We observe that
\begin{equation*}
\Phi _{i,j}\left( -x\right) =\frac{1}{4}\left( g_{i}\left( g_{j}\left(
-x+1\right) \right) -g_{i}\left( g_{j}\left( -x-1\right) \right) \right) =
\end{equation*}
\begin{equation*}
\frac{1}{4}\left( g_{i}\left( g_{j}\left( -\left( x-1\right) \right) \right)
-g_{i}\left( g_{j}\left( -\left( x+1\right) \right) \right) \right) =
\end{equation*}
\begin{equation*}
\frac{1}{4}\left( g_{i}\left( -g_{j}\left( x-1\right) \right) -g_{i}\left(
-g_{j}\left( x+1\right) \right) \right) =
\end{equation*}
\begin{equation*}
\frac{1}{4}\left( -g_{i}\left( g_{j}\left( x-1\right) \right) +g_{i}\left(
g_{j}\left( x+1\right) \right) \right) =
\end{equation*}
\begin{equation*}
\frac{1}{4}\left( g_{i}g_{j}\left( x+1\right) -g_{i}g_{j}\left( x-1\right)
\right) =\Phi _{i,j}\left( x\right) ,
\end{equation*}
that is
\begin{equation*}
\Phi _{i,j}\left( -x\right) =\Phi _{i,j}\left( x\right) ,\text{ \ }\forall
\text{ }x\in \mathbb{R}.
\end{equation*}
So $\Phi _{i,j}$ can serve as a density function in general.

So we have $g_{j}:\mathbb{R}\rightarrow \left( -1,1\right) $, $%
g_{i}|_{\left( -1,1\right) }:\left( -1,1\right) \rightarrow \left(
-1,1\right) $, and the strictly increasing function $H_{i,j}:=g_{i}|_{\left(
-1,1\right) }\circ g_{j}:\mathbb{R}\rightarrow \left( -1,1\right) $, with
the graph of $H_{i,j}$ containing an arc of finite length, such that $%
H_{i,j}\left( 0\right) =0$, starting at $\left( -1,g_{i}\left( g_{j}\left(
-1\right) \right) \right) $ and terminating at $\left( 1,g_{i}\left(
g_{j}\left( 1\right) \right) \right) .$We call this arc also $H_{i,j}$. In
particular $H_{i,j}$ is negative and convex over $(-1,0]$, and it is
positive and concave over $[0,1)$.

So it has compact support $\left[ -1,1\right] $ and it is like a squashing
function, see [4], Ch. 1, p. 8.

We will work from now on with $\left \vert H_{i,j}\right \vert $, which has
as a graph a cusp joining the points $\left( -1,\left \vert g_{i}\left(
g_{j}\left( -1\right) \right) \right \vert \right) $, $\left( 0,0\right)
,\left( 1,g_{i}\left( g_{j}\left( 1\right) \right) \right) $ and with
compact support, again, $\left[ -1,1\right] $. The points $\left(
-1,\left
\vert g_{i}\left( g_{j}\left( -1\right) \right) \right \vert
\right) $, $\left( 1,g_{i}\left( g_{j}\left( 1\right) \right) \right) $
belong to the graph of $\left \vert H_{i,j}\right \vert $ and $\left(
0,0\right) $ too.

Typically $H_{i,j}$ has a steeper slope than of $g_{j}$, but it is flatter
and closer to the $x$-axis than $g_{j}$ is, e.g. $\tanh \left( \tanh
x\right) $ has asymptotes $\pm 0.76$, while $\tanh x$ has asymptotes $\pm 1$
, notice that $\tanh \left( 1\right) =0.76$. Clearly $H_{i,j}$ has
applications in spiking neural networks.

Here we consider functions $f:\mathbb{R}\rightarrow \mathbb{R}$ to be either
continuous and bounded, or uniformly continuous.

The first modulus of continuity is given by
\begin{equation*}
\omega _{1}\left( f,\delta \right) :=\underset{%
\begin{array}{c}
x,y\in \mathbb{R}: \\
\left \vert x-y\right \vert \leq \delta%
\end{array}%
}{\sup }\left \vert f\left( x\right) -f\left( y\right) \right \vert ,\text{
\ }\delta >0.
\end{equation*}

Here we have that $\omega _{1}\left( f,\delta \right) <+\infty $, $\delta
>0. $

In this article we study the generalized fractional pointwise and uniform
convergences with rates over the real line, to the unit operator, of the 49
different "normalized cusp neural network operators",
\begin{equation}
\left( A_{n}^{\left( i,j\right) }\left( f\right) \right) \left( x\right) :=%
\frac{\sum \limits_{k=-n^{2}}^{n^{2}}f\left( \frac{k}{n}\right) \left \vert
H_{i,j}\left( n^{1-\alpha }\left( x-\frac{k}{n}\right) \right) \right \vert
}{\sum \limits_{k=-n^{2}}^{n^{2}}\left \vert H_{i,j}\left( n^{1-\alpha
}\left( x-\frac{k}{n}\right) \right) \right \vert },  \tag{9}  \label{9.}
\end{equation}
where $0<\alpha <1$ and $x\in \mathbb{R}$, $n\in \mathbb{N}$; $i,j\in
\left
\{ 1,...,7\right \} ,$ $i\neq j$ or $i=j.$

Notice $A_{n}^{\left( i,j\right) }$ is a positive linear operator with $%
A_{n}^{\left( i,j\right) }\left( 1\right) =1$.

The terms in the ratio of sums (\ref{9.}) can be nonnegative and make sense,
iff $0<\left \vert n^{1-\alpha }\left( x-\frac{k}{n}\right) \right \vert
\leq 1$, i.e. $0<\left \vert x-\frac{k}{n}\right \vert \leq \frac{1}{%
n^{1-\alpha }},$ iff
\begin{equation}
nx-n^{\alpha }\leq x\leq nx+n^{\alpha }\text{, \ }x\neq \frac{k}{n}.
\tag{10}  \label{10.}
\end{equation}
In order to have the desired order of numbers
\begin{equation}
-n^{2}\leq nx-n^{\alpha }\leq nx+n^{\alpha }\leq n^{2}\text{, \ }x\neq \frac{%
k}{n},  \tag{11}  \label{11.}
\end{equation}
it is sufficient to assume that
\begin{equation}
n\geq 1+\left \vert x\right \vert \text{, \ }x\neq \frac{k}{n}.  \tag{12}
\label{12.}
\end{equation}
When $x\in \left[ -1,1\right] $, $x\neq \frac{k}{n}$, it is enough to assume
$n\geq 2$, which implies (\ref{11.}), and $x\neq \frac{k}{n}.$

But the unique case $x=\frac{k}{n}$ contributes nothing and can be ignored.

Thus, without loss of generality we can take always that $x\neq \frac{k}{n}.$

\begin{proposition}
\label{p2}(\cite{2}) Let $a,b\in \mathbb{R}$, $a\leq b$. Let $card\left(
k\right) $ $\left( \geq 0\right) $ be the maximum number of integers
contained in $\left[ a,b\right] $. Then
\begin{equation*}
\max \left( 0,\left( b-a\right) -1\right) \leq card\left( k\right) \leq
\left( b-a\right) +1.
\end{equation*}
\end{proposition}

\begin{note}
\label{n3}We would like to establish a lower bound on $card\left( k\right) $
over the interval $\left[ nx-n^{\alpha },nx+n^{\alpha }\right] $. By
Proposition \ref{p2} we get that
\begin{equation*}
card\left( k\right) \geq \max \left( 2n^{\alpha }-1,0\right) .
\end{equation*}
We obtain $card\left( k\right) \geq 1$, if $2n^{\alpha }-1\geq 1$ iff $n\geq
1,$ which is always true.

So to have the desired order $\left( 11\right) $ and $card\left( k\right)
\geq 1$ over $\left[ nx-n^{\alpha },nx+n^{\alpha }\right] ,$ it is enough to
consider
\begin{equation}
n\geq \max \left( 1+\left \vert x\right \vert ,1\right) =1+\left \vert
x\right \vert .  \tag{13}  \label{13.}
\end{equation}
Also notice that $card\left( k\right) \rightarrow +\infty $, as $%
n\rightarrow +\infty .$
\end{note}

Denote by $\left[ \cdot \right] $ the integral part of a number and $%
\left
\lceil \cdot \right \rceil $ its ceiling.

Thus, it is clear that
\begin{equation}
\left( A_{n}^{\left( i,j\right) }\left( f\right) \right) \left( x\right) :=%
\frac{\sum \limits_{k=\left \lceil nx-n^{\alpha }\right \rceil }^{\left[
nx+n^{\alpha }\right] }f\left( \frac{k}{n}\right) \left \vert H_{i,j}\left(
n^{1-\alpha }\left( x-\frac{k}{n}\right) \right) \right \vert }{\sum
\limits_{k=\left \lceil nx-n^{\alpha }\right \rceil }^{\left[ nx+n^{\alpha }%
\right] }\left \vert H_{i,j}\left( n^{1-\alpha }\left( x-\frac{k}{n}\right)
\right) \right \vert },  \tag{14}  \label{14.}
\end{equation}
$0<\alpha <1$, and $n\in \mathbb{N}:n\geq 1+\left \vert x\right \vert ,$ $%
x\in \mathbb{R};$ all $i,j\in \left \{ 1,...,7\right \} ,$ $i\neq j$ or $%
i=j. $

\section{\; About Generalized Fractional Calculus}

We will use

\begin{definition}
\label{d4}(\cite{5}) Let $\beta >0$, $\left \lceil \beta \right \rceil =N$, $%
\left \lceil \cdot \right \rceil $ the ceiling of the number. Let $f\in
C^{N}\left( \left[ a,b\right] ,X\right) $, where $\left[ a,b\right] \subset
\mathbb{R}$, and $\left( X,\left \Vert \cdot \right \Vert \right) $ is a
Banach space. Let $g\in C^{1}\left( \left[ a,b\right] \right) $, strictly
increasing, such that $g^{-1}\in C^{N}\left( \left[ g\left( a\right)
,g\left( b\right) \right] \right) .$

We define the left generalized $g$-fractional derivative $X$-valued of $f$
of order $\beta $ as follows:
\begin{equation}
\left( D_{\ast a;g}^{\beta }f\right) \left( x\right) :=\frac{1}{\Gamma
\left( N-\beta \right) }\int_{a}^{x}\left( g\left( x\right) -g\left(
t\right) \right) ^{N-\beta -1}g^{\prime }\left( t\right) \left( f\circ
g^{-1}\right) ^{\left( N\right) }\left( g\left( t\right) \right) dt,
\tag{15}  \label{15}
\end{equation}
$\forall $ $x\in \left[ a,b\right] $. The last integral is of Bochner type.

If $\beta \notin \mathbb{N}$, by \cite{5}, we have that $\left( D_{\ast
a;g}^{\beta }f\right) \in C\left( \left[ a,b\right] ,X\right) .$
\end{definition}

We set
\begin{equation}
D_{\ast a;g}^{N}f\left( x\right) :=\left( \left( f\circ g^{-1}\right)
^{\left( N\right) }\circ g\right) \left( x\right) \in C\left( \left[ a,b%
\right] ,X\right) \text{, \ }N\in \mathbb{N},  \tag{16}  \label{16}
\end{equation}
\begin{equation*}
D_{\ast a;g}^{0}f\left( x\right) =f\left( x\right) \text{, \ }\forall \text{
}x\in \left[ a,b\right] .
\end{equation*}
When $g=id$, then
\begin{equation}
D_{\ast a;g}^{\beta }f=D_{\ast a;id}^{\beta }f=D_{\ast a}^{\beta }f,
\tag{17}  \label{17}
\end{equation}
the usual left $X$-valued Caputo fractional derivative, see \cite{9,11}.

We will use

\begin{definition}
\label{d5}(\cite{5})Let $\beta >0,$ $\left \lceil \beta \right \rceil =N$, $%
\left \lceil \cdot \right \rceil $ the ceiling of the number. Let $f\in
C^{N}\left( \left[ a,b\right] ,X\right) $, where $\left[ a,b\right] \subset
\mathbb{R}$, and $\left( X,\left \Vert \cdot \right \Vert \right) $ is a
Banach space. Let $g\in C^{1}\left( \left[ a,b\right] \right) $, strictly
increasing, such that $g^{-1}\in C^{N}\left( \left[ g\left( a\right)
,g\left( b\right) \right] \right) .$

We define the right generalized $g$-fractional derivative $X$-valued of $f$
of order $\beta $ as follows:
\begin{equation}
\left( D_{b-;g}^{\beta }f\right) \left( x\right) :=\frac{\left( -1\right)
^{N}}{\Gamma \left( N-\beta \right) }\int_{x}^{b}\left( g\left( t\right)
-g\left( x\right) \right) ^{N-\beta -1}g^{\prime }\left( t\right) \left(
f\circ g^{-1}\right) ^{\left( N\right) }\left( g\left( t\right) \right) dt,
\tag{18}  \label{18}
\end{equation}
$\forall $ $x\in \left[ a,b\right] $. The last integral is of Bochner type.

If $\beta \notin \mathbb{N}$, by \cite{5}, we have that $\left(
D_{b-;g}^{\beta }f\right) \in C\left( \left[ a,b\right] ,X\right) .$
\end{definition}

We set
\begin{equation}
D_{b-;g}^{N}f\left( x\right) :=\left( -1\right) ^{N}\left( \left( f\circ
g^{-1}\right) ^{\left( N\right) }\circ g\right) \left( x\right) \in C\left(
\left[ a,b\right] ,X\right) \text{, \ }N\in \mathbb{N},  \tag{19}  \label{19}
\end{equation}
\begin{equation*}
D_{b-;g}^{0}f\left( x\right) =f\left( x\right) \text{, \ }\forall \text{ }%
x\in \left[ a,b\right] .
\end{equation*}
When $g=id$, then
\begin{equation}
D_{b-;g}^{\beta }f=D_{b-;id}^{\beta }f=D_{b-}^{\beta }f,  \tag{20}
\label{20}
\end{equation}
the usual right $X$-valued Caputo fractional derivative, see \cite{3,10,11}.

We mention the following $g$-left generalized $X$-valued Taylor's formula:

\begin{theorem}
\label{t6}(\cite{5}) Let $\beta >0,$ $N=\left\lceil \beta \right\rceil $ and
$f\in C^{N}\left( \left[ a,b\right] ,X\right) $, where $\left[ a,b\right]
\subset \mathbb{R}$, and $\left( X,\left\Vert \cdot \right\Vert \right) $ is
a Banach space. Let $g\in C^{1}\left( \left[ a,b\right] \right) $, strictly
increasing, such that $g^{-1}\in C^{N}\left( \left[ g\left( a\right)
,g\left( b\right) \right] \right) .$ Then
\begin{equation*}
f\left( x\right) =f\left( a\right) +\sum_{j_{\ast }=1}^{N-1}\frac{\left(
g\left( x\right) -g\left( a\right) \right) ^{j_{\ast }}}{j_{\ast }!}\left(
f\circ g^{-1}\right) ^{\left( j_{\ast }\right) }\left( g\left( a\right)
\right) +
\end{equation*}
\begin{equation*}
\frac{1}{\Gamma \left( \beta \right) }\int_{a}^{x}\left( g\left( x\right)
-g\left( t\right) \right) ^{\beta -1}g^{\prime }\left( t\right) \left(
D_{\ast a;g}^{\beta }f\right) \left( t\right) dt=
\end{equation*}
\begin{equation}
f\left( a\right) +\sum_{j_{\ast }=1}^{N-1}\frac{\left( g\left( x\right)
-g\left( a\right) \right) ^{j_{\ast }}}{j_{\ast }!}\left( f\circ
g^{-1}\right) ^{\left( j_{\ast }\right) }\left( g\left( a\right) \right) +
\tag{21}  \label{21}
\end{equation}
\begin{equation*}
\frac{1}{\Gamma \left( \beta \right) }\int_{g\left( a\right) }^{g\left(
x\right) }\left( g\left( x\right) -z\right) ^{\beta -1}\left( \left( D_{\ast
a;g}^{\beta }f\right) \circ g^{-1}\right) \left( z\right) dz,\text{ \ }%
\forall \text{ }x\in \left[ a,b\right] .
\end{equation*}
\end{theorem}

We also mention the following $g$-right generalized $X$-valued Taylor's
formula:

\begin{theorem}
\label{t7}(\cite{5}) Let $\beta >0,$ $N=\left \lceil \beta \right \rceil $
and $f\in C^{N}\left( \left[ a,b\right] ,X\right) $, where $\left[ a,b\right]
\subset \mathbb{R}$, and $\left( X,\left \Vert \cdot \right \Vert \right) $
is a Banach space. Let $g\in C^{1}\left( \left[ a,b\right] \right) $,
strictly increasing, such that $g^{-1}\in C^{N}\left( \left[ g\left(
a\right) ,g\left( b\right) \right] \right) .$ Then
\begin{equation*}
f\left( x\right) =f\left( b\right) +\sum_{j_{\ast }=1}^{N-1}\frac{\left(
g\left( x\right) -g\left( b\right) \right) ^{j_{\ast }}}{j_{\ast }!}\left(
f\circ g^{-1}\right) ^{\left( j_{\ast }\right) }\left( g\left( b\right)
\right) +
\end{equation*}
\begin{equation*}
\frac{1}{\Gamma \left( \beta \right) }\int_{x}^{b}\left( g\left( t\right)
-g\left( x\right) \right) ^{\beta -1}g^{\prime }\left( t\right) \left(
D_{b-;g}^{\beta }f\right) \left( t\right) dt=
\end{equation*}
\begin{equation}
f\left( b\right) +\sum_{j_{\ast }=1}^{N-1}\frac{\left( g\left( x\right)
-g\left( b\right) \right) ^{j_{\ast }}}{j_{\ast }!}\left( f\circ
g^{-1}\right) ^{\left( j_{\ast }\right) }\left( g\left( b\right) \right) +
\tag{22}  \label{22}
\end{equation}
\begin{equation*}
\frac{1}{\Gamma \left( \beta \right) }\int_{g\left( x\right) }^{g\left(
b\right) }\left( z-g\left( x\right) \right) ^{\beta -1}\left( \left(
D_{b-;g}^{\beta }f\right) \circ g^{-1}\right) \left( z\right) dz,\text{ \ }%
\forall \text{ }x\in \left[ a,b\right] .
\end{equation*}
\end{theorem}

Integrals in (\ref{15})-(\ref{22}) are of Bochner type. For the Bochner
integral excellent resources are \cite{17} and \cite[pp. 422-428]{1}.

When $X=\mathbb{R}$, then the Bochner integral collapses to the Lebesgue
integral. So all the above apply also to real valued functions which we
treat here.

\begin{example}
\label{e8} The following function
\begin{equation*}
f\left( x\right) =x+\sqrt{x^{2}+1},\text{ \ }x\in \mathbb{R}\text{,}
\end{equation*}
has derivative $f^{\prime }\left( x\right) =1+\frac{x}{\sqrt{x^{2}+1}}$ with
range $(0,2)$.

Indeed we have: $1+\frac{x}{\sqrt{x^{2}+1}}>0$, iff $\frac{x}{\sqrt{x^{2}+1}}%
>-1$, iff $x>-\sqrt{x^{2}+1}$, iff $-x<\sqrt{x^{2}+1}$. The last is true if $%
x\geq 0,$ and if $x<0$, then $\left \vert x\right \vert <\sqrt{x^{2}+1}$,
iff $x^{2}<x^{2}+1$, iff $0<1$ true.

And we have $1+\frac{x}{\sqrt{x^{2}+1}}<2$, iff $\frac{x}{\sqrt{x^{2}+1}}<1$
, iff $x<\sqrt{x^{2}+1}$, if $x<0$ the last is true, if $x>0$ then $%
x^{2}<x^{2}+1$, iff $0<1$ true. Thus, $f$ is strictly increasing,
differentiable and uniformly continuous. It has $C^{1}$ inverse, i.e. $%
f^{-1} $ is continuous and differentiable, with $\left( f^{-1}\right)
^{\prime }\left( y\right) =\frac{1}{f^{\prime }\left( f^{-1}\left( y\right)
\right) }$, by continuous $f^{\prime }$ and $f^{\prime }\neq 0.$
\end{example}

\begin{example}
\label{e9} Let $f\left( x\right) =x+\arctan \left( x\right) $. Then $%
f^{\prime }\left( x\right) =1+\frac{1}{1+x^{2}}$, with $1<f^{\prime }\left(
x\right) \leq 2$, $\forall $ $x\in \mathbb{R}$. Thus $f$ is $C^{1}$,
strictly increasing and uniformly continuous.

By the inverse function theorem $f^{-1}$ is differentiable and $\left(
f^{-1}\right) ^{\prime }$ is continuous.
\end{example}

\begin{example}
\label{e10} Let $f\left( x\right) =2x+\sin \left( x\right) $, it is $1\leq
f^{\prime }\left( x\right) \leq 3$, thus $f$ is uniformly continuous,
strictly increasing, continuously differentiable with a $C^{1}$ inverse.
\end{example}

So our assumption on $g$ in Definitions \ref{d4}, \ref{d5} and Theorems \ref%
{t6}, \ref{t7} is possible.

\section{\; Main Results}

We present the following generalized fractional approximation result.

\begin{theorem}
\label{t11} Let $\beta >0$, $\beta \notin \mathbb{N}$, $N=\left \lceil \beta
\right \rceil $. Let also a strictly increasing $g\in C^{1}\left( \mathbb{R}%
\right) $, such that $g^{-1}\in C^{N}\left( \mathbb{R}\right) $; $g$ is also
assumed to be uniformly continuous. Let also $f\in C^{N}\left( \mathbb{R}%
\right) $ and assume that $\left( f\circ g^{-1}\right) ^{\left( N\right)
}\circ g\in L_{\infty }\left( \mathbb{R}\right) $. Let also $x\in \mathbb{R}$%
, $n\in \mathbb{N}:n\geq 1+\left \vert x\right \vert $. We further assume
that $D_{\ast x;g}^{\beta }f$, $D_{x-;g}^{\beta }f$ are uniformly continuous
or continuous and bounded on $[x,+\infty )$, $(-\infty ,x]$, respectively.
Then

1)
\begin{equation}
\underset{i\neq j\text{ or }i=j}{\underset{i,j\in \left \{ 1,...,7\right \} ,%
}{\max }}\left \vert A_{n}^{\left( i,j\right) }\left( f\right) \left(
x\right) -f\left( x\right) -\sum_{j_{\ast }=1}^{N-1}\frac{\left( f\circ
g^{-1}\right) ^{\left( j_{\ast }\right) }\left( g\left( x\right) \right) }{%
j_{\ast }!}\right.  \tag{23}  \label{23}
\end{equation}
\begin{equation*}
\left. \left( \frac{\sum \limits_{k=\left \lceil nx-n^{\alpha }\right \rceil
}^{\left[ nx+n^{\alpha }\right] }\left( g\left( \frac{k}{n}\right) -g\left(
x\right) \right) ^{j\ast }\left \vert H_{i,j}\left( n^{1-\alpha }\left( x-%
\frac{k}{n}\right) \right) \right \vert }{\sum \limits_{k=\left \lceil
nx-n^{\alpha }\right \rceil }^{\left[ nx+n^{\alpha }\right] }\left \vert
H_{i,j}\left( n^{1-\alpha }\left( x-\frac{k}{n}\right) \right) \right \vert }%
\right) \right \vert \leq
\end{equation*}
\begin{equation*}
\frac{\omega _{1}^{\beta }\left( g,\frac{1}{n^{1-\alpha }}\right) }{\Gamma
\left( \beta +1\right) }\left[ \omega _{1}\left( D_{x-;g}^{\beta }f;\frac{1}{%
n^{1-\alpha }}\right) _{(-\infty ,x]}+\omega _{1}\left( D_{\ast x;g}^{\beta
}f;\frac{1}{n^{1-\alpha }}\right) _{[x,+\infty )}\right] ,
\end{equation*}
above it is $\sum_{j=1}^{0}\cdot =0.$

2) If $N=1$ or $\left( f\circ g^{-1}\right) ^{\left( j\right) }\left(
g\left( x\right) \right) =0$, $j=1,...,N-1$, we have
\begin{equation}
\underset{i\neq j\text{ or }i=j}{\underset{i,j\in \left\{ 1,...,7\right\} ,}{%
\max }}\left\vert A_{n}^{\left( i,j\right) }\left( f\right) \left( x\right)
-f\left( x\right) \right\vert \leq   \tag{24}  \label{24}
\end{equation}
\begin{equation*}
\frac{\omega _{1}^{\beta }\left( g,\frac{1}{n^{1-\alpha }}\right) }{\Gamma
\left( \beta +1\right) }\left[ \omega _{1}\left( D_{x-;g}^{\beta }f;\frac{1}{%
n^{1-\alpha }}\right) _{(-\infty ,x]}+\omega _{1}\left( D_{\ast x;g}^{\beta
}f;\frac{1}{n^{1-\alpha }}\right) _{[x,+\infty )}\right] ,
\end{equation*}

3)
\begin{equation}
\underset{i\neq j\text{ or }i=j}{\underset{i,j\in \left \{ 1,...,7\right \} ,%
}{\max }}\left \vert A_{n}^{\left( i,j\right) }\left( f\right) \left(
x\right) -f\left( x\right) \right \vert \leq \sum_{j_{\ast }=1}^{N-1}\frac{%
\left \vert \left( f\circ g^{-1}\right) ^{\left( j_{\ast }\right) }\left(
g\left( x\right) \right) \right \vert }{j_{\ast }!}\omega _{1}^{j_{\ast
}}\left( g,\frac{1}{n^{1-\alpha }}\right)  \tag{25}  \label{25}
\end{equation}
\begin{equation*}
+\frac{\omega _{1}^{\beta }\left( g,\frac{1}{n^{1-\alpha }}\right) }{\Gamma
\left( \beta +1\right) }\left[ \omega _{1}\left( D_{x-;g}^{\beta }f;\frac{1}{%
n^{1-\alpha }}\right) _{(-\infty ,x]}+\omega _{1}\left( D_{\ast x;g}^{\beta
}f;\frac{1}{n^{1-\alpha }}\right) _{[x,+\infty )}\right] .
\end{equation*}
Here we get generalized fractionally with rates the pointwise convergence of
$A_{n}^{\left( i,j\right) }\left( f\right) \left( x\right) \rightarrow
f\left( x\right) $, as $n\rightarrow +\infty $, $x\in \mathbb{R}$; $i,j\in
\left \{ 1,...,7\right \} .$
\end{theorem}

\begin{proof}
Let $x\in \mathbb{R}$. We have that $D_{x-;g}^{\beta }f\left( x\right)
=D_{\ast x;g}^{\beta }f\left( x\right) =0$, explanation at the end of this
proof.

From Theorem 6, by using the left generalized Caputo fractional Taylor
formula we get that
\begin{equation*}
f\left( \frac{k}{n}\right) =\sum_{j_{\ast }=0}^{N-1}\frac{\left( f\circ
g^{-1}\right) ^{\left( j_{\ast }\right) }\left( g\left( x\right) \right) }{%
j_{\ast }!}\left( g\left( \frac{k}{n}\right) -g\left( x\right) \right)
^{j_{\ast }}+
\end{equation*}
\begin{equation}
\frac{1}{\Gamma \left( \beta \right) }\int_{x}^{\frac{k}{n}}\left( g\left(
\frac{k}{n}\right) -g\left( J\right) \right) ^{\beta -1}g^{\prime }\left(
J\right) \left[ \left( D_{\ast x;g}^{\beta }f\right) \left( J\right) -\left(
D_{\ast x;g}^{\beta }f\right) \left( x\right) \right] dJ,  \tag{26}
\label{26}
\end{equation}
for all $x\leq \frac{k}{n}\leq x+n^{\alpha -1}$, iff $\left\lceil
nx\right\rceil \leq k\leq \left[ nx+n^{\alpha }\right] $, where $k\in
\mathbb{Z}.$

Also from Theorem 7, by using the right generalized Caputo fractional Taylor
formula we get
\begin{equation*}
f\left( \frac{k}{n}\right) =\sum_{j_{\ast }=0}^{N-1}\frac{\left( f\circ
g^{-1}\right) ^{\left( j_{\ast }\right) }\left( g\left( x\right) \right) }{%
j_{\ast }!}\left( g\left( \frac{k}{n}\right) -g\left( x\right) \right)
^{j_{\ast }}+
\end{equation*}
\begin{equation}
\frac{1}{\Gamma \left( \beta \right) }\int_{\frac{k}{n}}^{x}\left( g\left(
J\right) -g\left( \frac{k}{n}\right) \right) ^{\beta -1}g^{\prime }\left(
J\right) \left[ \left( D_{x-;g}^{\beta }f\right) \left( J\right) -\left(
D_{x-;g}^{\beta }f\right) \left( x\right) \right] dJ,  \tag{27}  \label{27}
\end{equation}
for all $x-n^{\alpha -1}\leq \frac{k}{n}\leq x$, iff $\left\lceil
nx-n^{\alpha }\right\rceil \leq k\leq \left[ nx\right] $, where $k\in
\mathbb{Z}.$

Notice that $\left \lceil nx\right \rceil \leq \left[ nx\right] +1.$

Call
\begin{equation*}
W_{i,j}\left( x\right) :=\sum \limits_{k=\left \lceil nx-n^{\alpha }\right
\rceil }^{\left[ nx+n^{\alpha }\right] }\left \vert H_{i,j}\left(
n^{1-\alpha }\left( x-\frac{k}{n}\right) \right) \right \vert ,\text{ \ }%
i,j\in \left \{ 1,...,7\right \} .
\end{equation*}
Hence we have
\begin{equation*}
\frac{f\left( \frac{k}{n}\right) \left \vert H_{i,j}\left( n^{1-\alpha
}\left( x-\frac{k}{n}\right) \right) \right \vert }{W_{i,j}\left( x\right) }=
\end{equation*}
\begin{equation}
\sum_{j_{\ast }=0}^{N-1}\frac{\left( f\circ g^{-1}\right) ^{\left( j_{\ast
}\right) }\left( g\left( x\right) \right) }{j_{\ast }!}\left( g\left( \frac{k%
}{n}\right) -g\left( x\right) \right) ^{j\ast }\frac{\left \vert
H_{i,j}\left( n^{1-\alpha }\left( x-\frac{k}{n}\right) \right) \right \vert
}{W_{i,j}\left( x\right) }+  \tag{28}  \label{28}
\end{equation}
\begin{equation*}
\frac{\left \vert H_{i,j}\left( n^{1-\alpha }\left( x-\frac{k}{n}\right)
\right) \right \vert }{W_{i,j}\left( x\right) \Gamma \left( \beta \right) }%
\int_{x}^{\frac{k}{n}}\left( g\left( \frac{k}{n}\right) -g\left( J\right)
\right) ^{\beta -1}g^{\prime }\left( J\right) \left[ D_{\ast x;g}^{\beta
}f\left( J\right) -D_{\ast x;g}^{\beta }f\left( x\right) \right] dJ,
\end{equation*}
and
\begin{equation*}
\frac{f\left( \frac{k}{n}\right) \left \vert H_{i,j}\left( n^{1-\alpha
}\left( x-\frac{k}{n}\right) \right) \right \vert }{W_{i,j}\left( x\right) }=
\end{equation*}
\begin{equation}
\sum_{j_{\ast }=0}^{N-1}\frac{\left( f\circ g^{-1}\right) ^{\left( j_{\ast
}\right) }\left( g\left( x\right) \right) }{j_{\ast }!}\left( g\left( \frac{k%
}{n}\right) -g\left( x\right) \right) ^{j\ast }\frac{\left \vert
H_{i,j}\left( n^{1-\alpha }\left( x-\frac{k}{n}\right) \right) \right \vert
}{W_{i,j}\left( x\right) }+  \tag{29}  \label{29}
\end{equation}
\begin{equation*}
\frac{\left \vert H_{i,j}\left( n^{1-\alpha }\left( x-\frac{k}{n}\right)
\right) \right \vert }{W_{i,j}\left( x\right) \Gamma \left( \beta \right) }%
\int_{\frac{k}{n}}^{x}\left( g\left( J\right) -g\left( \frac{k}{n}\right)
\right) ^{\beta -1}g^{\prime }\left( J\right) \left[ D_{x-;g}^{\beta
}f\left( J\right) -D_{x-;g}^{\beta }f\left( x\right) \right] dJ.
\end{equation*}
Therefore we obtain
\begin{equation*}
\frac{\sum_{k=\left[ nx\right] +1}^{\left[ nx+n^{\alpha }\right] }f\left(
\frac{k}{n}\right) \left \vert H_{i,j}\left( n^{1-\alpha }\left( x-\frac{k}{n%
}\right) \right) \right \vert }{W_{i,j}\left( x\right) }=
\end{equation*}
\begin{equation*}
\sum_{j_{\ast }=0}^{N-1}\frac{\left( f\circ g^{-1}\right) ^{\left( j_{\ast
}\right) }\left( g\left( x\right) \right) }{j_{\ast }!}\left( \frac{\sum_{k=%
\left[ nx\right] +1}^{\left[ nx+n^{\alpha }\right] }\left( g\left( \frac{k}{n%
}\right) -g\left( x\right) \right) ^{j\ast }\left \vert H_{i,j}\left(
n^{1-\alpha }\left( x-\frac{k}{n}\right) \right) \right \vert }{%
W_{i,j}\left( x\right) }\right)
\end{equation*}
\begin{equation*}
+\sum_{k=\left[ nx\right] +1}^{\left[ nx+n^{\alpha }\right] }\frac{\left
\vert H_{i,j}\left( n^{1-\alpha }\left( x-\frac{k}{n}\right) \right) \right
\vert }{W_{i,j}\left( x\right) \Gamma \left( \beta \right) }
\end{equation*}
\begin{equation}
\int_{x}^{\frac{k}{n}}\left( g\left( \frac{k}{n}\right) -g\left( J\right)
\right) ^{\beta -1}g^{\prime }\left( J\right) \left[ D_{\ast x;g}^{\beta
}f\left( J\right) -D_{\ast x;g}^{\beta }f\left( x\right) \right] dJ,
\tag{30}  \label{30}
\end{equation}
and
\begin{equation*}
\frac{\sum_{k=\left \lceil nx-n^{\alpha }\right \rceil }^{\left[ nx\right]
}f\left( \frac{k}{n}\right) \left \vert H_{i,j}\left( n^{1-\alpha }\left( x-%
\frac{k}{n}\right) \right) \right \vert }{W_{i,j}\left( x\right) }=
\end{equation*}
\begin{equation*}
\sum_{j_{\ast }=0}^{N-1}\frac{\left( f\circ g^{-1}\right) ^{\left( j_{\ast
}\right) }\left( g\left( x\right) \right) }{j_{\ast }!}\left( \frac{%
\sum_{k=\left \lceil nx-n^{\alpha }\right \rceil }^{\left[ nx\right] }\left(
g\left( \frac{k}{n}\right) -g\left( x\right) \right) ^{j\ast }\left \vert
H_{i,j}\left( n^{1-\alpha }\left( x-\frac{k}{n}\right) \right) \right \vert
}{W_{i,j}\left( x\right) }\right)
\end{equation*}
\begin{equation*}
+\sum_{k=\left \lceil nx-n^{\alpha }\right \rceil }^{\left[ nx\right] }\frac{%
\left \vert H_{i,j}\left( n^{1-\alpha }\left( x-\frac{k}{n}\right) \right)
\right \vert }{W_{i,j}\left( x\right) \Gamma \left( \beta \right) }
\end{equation*}
\begin{equation}
\int_{\frac{k}{n}}^{x}\left( g\left( J\right) -g\left( \frac{k}{n}\right)
\right) ^{\beta -1}g^{\prime }\left( J\right) \left[ D_{x-;g}^{\beta
}f\left( J\right) -D_{x-;g}^{\beta }f\left( x\right) \right] dJ.  \tag{31}
\label{31}
\end{equation}
Adding the two equalities (\ref{30}) and (\ref{31}) we obtain
\begin{equation}
\left( A_{n}^{\left( i,j\right) }\left( f\right) \right) \left( x\right)
=\sum_{j_{\ast }=0}^{N-1}\frac{\left( f\circ g^{-1}\right) ^{\left( j_{\ast
}\right) }\left( g\left( x\right) \right) }{j_{\ast }!}  \tag{32}  \label{32}
\end{equation}
\begin{equation*}
\left( \frac{\sum \limits_{k=\left \lceil nx-n^{\alpha }\right \rceil }^{%
\left[ nx+n^{\alpha }\right] }\left( g\left( \frac{k}{n}\right) -g\left(
x\right) \right) ^{j\ast }\left \vert H_{i,j}\left( n^{1-\alpha }\left( x-%
\frac{k}{n}\right) \right) \right \vert }{W_{i,j}\left( x\right) }\right) +%
\widehat{M}_{n}^{\left( i,j\right) }\left( x\right) ,
\end{equation*}
where $i,j\in \left \{ 1,...,7\right \} ,$ and
\begin{equation}
\widehat{M}_{n}^{\left( i,j\right) }\left( x\right) :=\sum_{k=\left \lceil
nx-n^{\alpha }\right \rceil }^{\left[ nx\right] }\frac{\left \vert
H_{i,j}\left( n^{1-\alpha }\left( x-\frac{k}{n}\right) \right) \right \vert
}{W_{i,j}\left( x\right) \Gamma \left( \beta \right) }  \tag{33}  \label{33}
\end{equation}
\begin{equation*}
\int_{\frac{k}{n}}^{x}\left( g\left( J\right) -g\left( \frac{k}{n}\right)
\right) ^{\beta -1}g^{\prime }\left( J\right) \left[ D_{x-;g}^{\beta
}f\left( J\right) -D_{x-;g}^{\beta }f\left( x\right) \right] dJ+
\end{equation*}
\begin{equation*}
\sum_{k=\left[ nx\right] +1}^{\left[ nx+n^{\alpha }\right] }\frac{\left
\vert H_{i,j}\left( n^{1-\alpha }\left( x-\frac{k}{n}\right) \right) \right
\vert }{W_{i,j}\left( x\right) \Gamma \left( \beta \right) }
\end{equation*}
\begin{equation*}
\int_{x}^{\frac{k}{n}}\left( g\left( \frac{k}{n}\right) -g\left( J\right)
\right) ^{\beta -1}g^{\prime }\left( J\right) \left[ D_{\ast x;g}^{\beta
}f\left( J\right) -D_{\ast x;g}^{\beta }f\left( x\right) \right] dJ.
\end{equation*}

We call
\begin{equation}
\widehat{M}_{1n}^{\left( i,j\right) }\left( x\right) :=\sum_{k=\left \lceil
nx-n^{\alpha }\right \rceil }^{\left[ nx\right] }\frac{\left \vert
H_{i,j}\left( n^{1-\alpha }\left( x-\frac{k}{n}\right) \right) \right \vert
}{W_{i,j}\left( x\right) \Gamma \left( \beta \right) }  \tag{34}  \label{34}
\end{equation}
\begin{equation*}
\int_{\frac{k}{n}}^{x}\left( g\left( J\right) -g\left( \frac{k}{n}\right)
\right) ^{\beta -1}g^{\prime }\left( J\right) \left[ D_{x-;g}^{\beta
}f\left( J\right) -D_{x-;g}^{\beta }f\left( x\right) \right] dJ,
\end{equation*}
and
\begin{equation}
\widehat{M}_{2n}^{\left( i,j\right) }\left( x\right) :=\sum_{k=\left[ nx%
\right] +1}^{\left[ nx+n^{\alpha }\right] }\frac{\left \vert H_{i,j}\left(
n^{1-\alpha }\left( x-\frac{k}{n}\right) \right) \right \vert }{%
W_{i,j}\left( x\right) \Gamma \left( \beta \right) }  \tag{35}  \label{35}
\end{equation}
\begin{equation*}
\int_{x}^{\frac{k}{n}}\left( g\left( \frac{k}{n}\right) -g\left( J\right)
\right) ^{\beta -1}g^{\prime }\left( J\right) \left[ D_{\ast x;g}^{\beta
}f\left( J\right) -D_{\ast x;g}^{\beta }f\left( x\right) \right] dJ.
\end{equation*}
That is
\begin{equation}
\widehat{M}_{n}^{\left( i,j\right) }\left( x\right) =\widehat{M}%
_{1n}^{\left( i,j\right) }\left( x\right) +\widehat{M}_{2n}^{\left(
i,j\right) }\left( x\right) ,  \tag{36}  \label{36}
\end{equation}
where $i,j\in \left \{ 1,...,7\right \} .$

Furthermore we get
\begin{equation*}
E_{n}^{\left( i,j\right) }\left( f\right) \left( x\right) :=\left(
A_{n}^{\left( i,j\right) }\left( f\right) \right) \left( x\right) -f\left(
x\right) -\sum_{j_{\ast }=1}^{N-1}\frac{\left( f\circ g^{-1}\right) ^{\left(
j_{\ast }\right) }\left( g\left( x\right) \right) }{j_{\ast }!}
\end{equation*}
\begin{equation}
\left( \frac{\sum \limits_{k=\left \lceil nx-n^{\alpha }\right \rceil }^{%
\left[ nx+n^{\alpha }\right] }\left( g\left( \frac{k}{n}\right) -g\left(
x\right) \right) ^{j\ast }\left \vert H_{i,j}\left( n^{1-\alpha }\left( x-%
\frac{k}{n}\right) \right) \right \vert }{W_{i,j}\left( x\right) }\right) =%
\widehat{M}_{n}^{\left( i,j\right) }\left( x\right) ,  \tag{37}  \label{37}
\end{equation}
that is
\begin{equation}
\left \vert E_{n}^{\left( i,j\right) }\left( f\right) \left( x\right) \right
\vert \leq \left \vert \widehat{M}_{n}^{\left( i,j\right) }\left( x\right)
\right \vert \leq \left \vert \widehat{M}_{1n}^{\left( i,j\right) }\left(
x\right) \right \vert +\left \vert \widehat{M}_{2n}^{\left( i,j\right)
}\left( x\right) \right \vert .  \tag{38}  \label{38}
\end{equation}

If $N=1$, then
\begin{equation}
E_{n}^{\left( i,j\right) }\left( f\right) \left( x\right) =A_{n}^{\left(
i,j\right) }\left( f\right) \left( x\right) -f\left( x\right) =\widehat{M}%
_{n}^{\left( i,j\right) }\left( x\right) .  \tag{39}  \label{39}
\end{equation}

If $\left( f\circ g^{-1}\right) ^{j_{\ast }}\left( g\left( x\right) \right)
=0$, for $j_{\ast }=1,...,N-1$, then again
\begin{equation}
E_{n}^{\left( i,j\right) }\left( f\right) \left( x\right) =A_{n}^{\left(
i,j\right) }\left( f\right) \left( x\right) -f\left( x\right) =\widehat{M}%
_{n}^{\left( i,j\right) }\left( x\right) ,  \tag{40}  \label{40}
\end{equation}
where $i,j\in \left \{ 1,...,7\right \} .$

We can also have
\begin{equation}
\left \vert A_{n}^{\left( i,j\right) }\left( f\right) \left( x\right)
-f\left( x\right) \right \vert \leq \sum_{j_{\ast }=1}^{N-1}\frac{\left
\vert \left( f\circ g^{-1}\right) ^{\left( j_{\ast }\right) }\left( g\left(
x\right) \right) \right \vert }{j_{\ast }!}\omega _{1}^{j_{\ast }}\left( g,%
\frac{1}{n^{1-\alpha }}\right) +\widehat{M}_{n}^{\left( i,j\right) }\left(
x\right) ,  \tag{41}  \label{41}
\end{equation}
by $g$ being also uniformly continuous we have that $\omega _{1}\left( g,%
\frac{1}{n^{1-\alpha }}\right) <+\infty .$

We observe that
\begin{equation*}
\left \vert \widehat{M}_{1n}^{\left( i,j\right) }\left( x\right) \right
\vert \leq \sum_{k=\left \lceil nx-n^{\alpha }\right \rceil }^{\left[ nx%
\right] }\frac{\left \vert H_{i,j}\left( n^{1-\alpha }\left( x-\frac{k}{n}%
\right) \right) \right \vert }{W_{i,j}\left( x\right) \Gamma \left( \beta
\right) }
\end{equation*}
\begin{equation}
\int_{\frac{k}{n}}^{x}\left( g\left( J\right) -g\left( \frac{k}{n}\right)
\right) ^{\beta -1}g^{\prime }\left( J\right) \left \vert D_{x-;g}^{\beta
}f\left( J\right) -D_{x-;g}^{\beta }f\left( x\right) \right \vert dJ\leq
\tag{42}  \label{42}
\end{equation}
\begin{equation*}
\sum_{k=\left \lceil nx-n^{\alpha }\right \rceil }^{\left[ nx\right] }\frac{%
\left \vert H_{i,j}\left( n^{1-\alpha }\left( x-\frac{k}{n}\right) \right)
\right \vert }{W_{i,j}\left( x\right) \Gamma \left( \beta \right) }
\end{equation*}
\begin{equation*}
\int_{\frac{k}{n}}^{x}\left( g\left( J\right) -g\left( \frac{k}{n}\right)
\right) ^{\beta -1}g^{\prime }\left( J\right) \omega _{1}\left(
D_{x-;g}^{\beta }f,\left \vert J-x\right \vert \right) _{(-\infty ,x]}dJ\leq
\end{equation*}
\begin{equation*}
\sum_{k=\left \lceil nx-n^{\alpha }\right \rceil }^{\left[ nx\right] }\frac{%
\left \vert H_{i,j}\left( n^{1-\alpha }\left( x-\frac{k}{n}\right) \right)
\right \vert }{W_{i,j}\left( x\right) \Gamma \left( \beta \right) }
\end{equation*}
\begin{equation*}
\left( \int_{g\left( \frac{k}{n}\right) }^{g\left( x\right) }\left( g\left(
J\right) -g\left( \frac{k}{n}\right) \right) ^{\beta -1}dg\left( J\right)
\right) \omega _{1}\left( D_{x-;g}^{\beta }f,\left \vert x-\frac{k}{n}\right
\vert \right) _{(-\infty ,x]}\leq
\end{equation*}
\begin{equation}
\left( \sum_{k=\left \lceil nx-n^{\alpha }\right \rceil }^{\left[ nx\right] }%
\frac{\left \vert H_{i,j}\left( n^{1-\alpha }\left( x-\frac{k}{n}\right)
\right) \right \vert }{W_{i,j}\left( x\right) \Gamma \left( \beta \right) }%
\frac{\left( g\left( x\right) -g\left( \frac{k}{n}\right) \right) ^{\beta }}{%
\beta }\right) \omega _{1}\left( D_{x-;g}^{\beta }f,\frac{1}{n^{1-\alpha }}%
\right) _{(-\infty ,x]}  \tag{43}  \label{43}
\end{equation}
(by assuming $g$ is uniformly continuous, then $\omega _{1}\left( g,\delta
\right) <+\infty ,$ $\forall $ $\delta >0$)
\begin{equation*}
\leq \left( \sum_{k=\left \lceil nx-n^{\alpha }\right \rceil }^{\left[ nx%
\right] }\frac{\left \vert H_{i,j}\left( n^{1-\alpha }\left( x-\frac{k}{n}%
\right) \right) \right \vert }{W_{i,j}\left( x\right) }\right) \frac{\omega
_{1}^{\beta }\left( g,\frac{1}{n^{1-\alpha }}\right) }{\Gamma \left( \beta
+1\right) }\omega _{1}\left( D_{x-;g}^{\beta }f,\frac{1}{n^{1-\alpha }}%
\right) _{(-\infty ,x]}
\end{equation*}
\begin{equation}
\leq \left( \frac{\sum_{k=\left \lceil nx-n^{\alpha }\right \rceil }^{\left[
nx+n^{\alpha }\right] }\left \vert H_{i,j}\left( n^{1-\alpha }\left( x-\frac{%
k}{n}\right) \right) \right \vert }{W_{i,j}\left( x\right) }\right) \frac{%
\omega _{1}^{\beta }\left( g,\frac{1}{n^{1-\alpha }}\right) }{\Gamma \left(
\beta +1\right) }\omega _{1}\left( D_{x-;g}^{\beta }f,\frac{1}{n^{1-\alpha }}%
\right) _{(-\infty ,x]}  \tag{44}  \label{44}
\end{equation}
\begin{equation*}
=\frac{\omega _{1}^{\beta }\left( g,\frac{1}{n^{1-\alpha }}\right) }{\Gamma
\left( \beta +1\right) }\omega _{1}\left( D_{x-;g}^{\beta }f,\frac{1}{%
n^{1-\alpha }}\right) _{(-\infty ,x]}.
\end{equation*}
Thus it holds
\begin{equation}
\left \vert \widehat{M}_{1n}^{\left( i,j\right) }\left( x\right) \right
\vert \leq \frac{\omega _{1}^{\beta }\left( g,\frac{1}{n^{1-\alpha }}\right)
}{\Gamma \left( \beta +1\right) }\omega _{1}\left( D_{x-;g}^{\beta }f,\frac{1%
}{n^{1-\alpha }}\right) _{(-\infty ,x]}.  \tag{45}  \label{45}
\end{equation}

Next we estimate
\begin{equation*}
\left \vert \widehat{M}_{2n}^{\left( i,j\right) }\left( x\right) \right
\vert \leq \frac{\sum_{k=\left[ nx\right] +1}^{\left[ nx+n^{\alpha }\right]
}\left \vert H_{i,j}\left( n^{1-\alpha }\left( x-\frac{k}{n}\right) \right)
\right \vert }{W_{i,j}\left( x\right) \Gamma \left( \beta \right) }
\end{equation*}
\begin{equation}
\int_{x}^{\frac{k}{n}}\left( g\left( \frac{k}{n}\right) -g\left( J\right)
\right) ^{\beta -1}g^{\prime }\left( J\right) \left \vert D_{\ast
x;g}^{\beta }f\left( J\right) -D_{\ast x;g}^{\beta }f\left( x\right) \right
\vert dJ\leq  \tag{46}  \label{46}
\end{equation}
\begin{equation*}
\frac{\sum_{k=\left[ nx\right] +1}^{\left[ nx+n^{\alpha }\right] }\left
\vert H_{i,j}\left( n^{1-\alpha }\left( x-\frac{k}{n}\right) \right) \right
\vert }{W_{i,j}\left( x\right) \Gamma \left( \beta \right) }
\end{equation*}
\begin{equation*}
\int_{x}^{\frac{k}{n}}\left( g\left( \frac{k}{n}\right) -g\left( J\right)
\right) ^{\beta -1}g^{\prime }\left( J\right) \omega _{1}\left( D_{\ast
x;g}^{\beta }f,\left \vert J-x\right \vert \right) _{[x,+\infty )}dJ\leq
\end{equation*}
\begin{equation}
\frac{\sum_{k=\left[ nx\right] +1}^{\left[ nx+n^{\alpha }\right] }\left
\vert H_{i,j}\left( n^{1-\alpha }\left( x-\frac{k}{n}\right) \right) \right
\vert }{W_{i,j}\left( x\right) \Gamma \left( \beta \right) }  \tag{47}
\label{47}
\end{equation}
\begin{equation*}
\left( \int_{g\left( x\right) }^{g\left( \frac{k}{n}\right) }\left( g\left(
\frac{k}{n}\right) -g\left( J\right) \right) ^{\beta -1}dg\left( J\right)
\right) \omega _{1}\left( D_{\ast x;g}^{\beta }f,\left \vert \frac{k}{n}%
-x\right \vert \right) _{[x,+\infty )}\leq
\end{equation*}
\begin{equation}
\left( \frac{\sum_{k=\left[ nx\right] +1}^{\left[ nx+n^{\alpha }\right]
}\left \vert H_{i,j}\left( n^{1-\alpha }\left( x-\frac{k}{n}\right) \right)
\right \vert }{W_{i,j}\left( x\right) \Gamma \left( \beta \right) }\frac{%
\left( g\left( \frac{k}{n}\right) -g\left( x\right) \right) ^{\beta }}{\beta
}\right) \omega _{1}\left( D_{\ast x;g}^{\beta }f,\frac{1}{n^{1-\alpha }}%
\right) _{[x,+\infty )}\leq  \tag{48}  \label{48}
\end{equation}
\begin{equation*}
\left( \frac{\sum_{k=\left[ nx\right] +1}^{\left[ nx+n^{\alpha }\right]
}\left \vert H_{i,j}\left( n^{1-\alpha }\left( x-\frac{k}{n}\right) \right)
\right \vert }{W_{i,j}\left( x\right) }\right) \frac{\omega _{1}^{\beta
}\left( g,\frac{1}{n^{1-\alpha }}\right) }{\Gamma \left( \beta +1\right) }%
\omega _{1}\left( D_{\ast x;g}^{\beta }f,\frac{1}{n^{1-\alpha }}\right)
_{[x,+\infty )}\leq
\end{equation*}
\begin{equation}
\left( \frac{\sum_{k=\left \lceil nx-n^{\alpha }\right \rceil }^{\left[
nx+n^{\alpha }\right] }\left \vert H_{i,j}\left( n^{1-\alpha }\left( x-\frac{%
k}{n}\right) \right) \right \vert }{W_{i,j}\left( x\right) }\right) \frac{%
\omega _{1}^{\beta }\left( g,\frac{1}{n^{1-\alpha }}\right) }{\Gamma \left(
\beta +1\right) }\omega _{1}\left( D_{\ast x;g}^{\beta }f,\frac{1}{%
n^{1-\alpha }}\right) _{[x,+\infty )}  \tag{49}  \label{49}
\end{equation}
\begin{equation*}
=\frac{\omega _{1}^{\beta }\left( g,\frac{1}{n^{1-\alpha }}\right) }{\Gamma
\left( \beta +1\right) }\omega _{1}\left( D_{\ast x;g}^{\beta }f,\frac{1}{%
n^{1-\alpha }}\right) _{[x,+\infty )}.
\end{equation*}
Thus
\begin{equation}
\left \vert \widehat{M}_{2n}^{\left( i,j\right) }\left( x\right) \right
\vert \leq \frac{\omega _{1}^{\beta }\left( g,\frac{1}{n^{1-\alpha }}\right)
}{\Gamma \left( \beta +1\right) }\omega _{1}\left( D_{\ast x;g}^{\beta }f,%
\frac{1}{n^{1-\alpha }}\right) _{[x,+\infty )}.  \tag{50}  \label{50}
\end{equation}

We have found that
\begin{equation*}
\left \vert \widehat{M}_{n}^{\left( i,j\right) }\left( x\right) \right \vert
\leq \frac{\omega _{1}^{\beta }\left( g,\frac{1}{n^{1-\alpha }}\right) }{%
\Gamma \left( \beta +1\right) }
\end{equation*}
\begin{equation}
\left[ \omega _{1}\left( D_{x-;g}^{\beta }f,\frac{1}{n^{1-\alpha }}\right)
_{(-\infty ,x]}+\omega _{1}\left( D_{\ast x;g}^{\beta }f,\frac{1}{%
n^{1-\alpha }}\right) _{[x,+\infty )}\right] <+\infty ,  \tag{51}  \label{51}
\end{equation}
where $i,j\in \left \{ 1,...,7\right \} .$

Here we assumed that $\left( f\circ g^{-1}\right) ^{\left( N\right) }\circ
g\in L_{\infty }\left( \mathbb{R}\right) .$

Next we get
\begin{equation*}
\left \vert D_{\ast x;g}^{\beta }f\left( s\right) \right \vert \leq \frac{1}{%
\Gamma \left( N-\beta \right) }\int_{x}^{s}\left( g\left( s\right) -g\left(
t\right) \right) ^{N-\beta -1}g^{\prime }\left( t\right) \left \vert \left(
f\circ g^{-1}\right) ^{\left( N\right) }\left( g\left( t\right) \right)
\right \vert dt\leq
\end{equation*}
\begin{equation}
\frac{\left \Vert \left( f\circ g^{-1}\right) ^{\left( N\right) }\circ
g\right \Vert _{\infty ,\mathbb{R}}}{\Gamma \left( N-\beta \right) }%
\int_{g\left( x\right) }^{g\left( s\right) }\left( g\left( s\right) -g\left(
t\right) \right) ^{N-\beta -1}dg\left( t\right) =  \tag{52}  \label{52}
\end{equation}
\begin{equation*}
\frac{\left \Vert \left( f\circ g^{-1}\right) ^{\left( N\right) }\circ
g\right \Vert _{\infty ,\mathbb{R}}}{\Gamma \left( N-\beta +1\right) }\left(
g\left( s\right) -g\left( x\right) \right) ^{N-\beta }.
\end{equation*}
That is
\begin{equation}
\left \vert D_{\ast x;g}^{\beta }f\left( s\right) \right \vert \leq \frac{%
\left \Vert \left( f\circ g^{-1}\right) ^{\left( N\right) }\circ g\right
\Vert _{\infty ,\mathbb{R}}}{\Gamma \left( N-\beta +1\right) }\left( g\left(
s\right) -g\left( x\right) \right) ^{N-\beta },  \tag{53}  \label{53}
\end{equation}
$\forall $ $s\in \lbrack x,+\infty ).$

Clearly then
\begin{equation}
D_{\ast x;g}^{\beta }f\left( x\right) =0.  \tag{54}  \label{54}
\end{equation}

Acting similarly we have:
\begin{equation*}
\left \vert D_{x-;g}^{\beta }f\left( s\right) \right \vert \leq \frac{1}{%
\Gamma \left( N-\beta \right) }\int_{s}^{x}\left( g\left( t\right) -g\left(
s\right) \right) ^{N-\beta -1}g^{\prime }\left( t\right) \left \vert \left(
f\circ g^{-1}\right) ^{\left( N\right) }\left( g\left( t\right) \right)
\right \vert dt\leq
\end{equation*}
\begin{equation}
\frac{\left \Vert \left( f\circ g^{-1}\right) ^{\left( N\right) }\circ
g\right \Vert _{\infty ,\mathbb{R}}}{\Gamma \left( N-\beta \right) }%
\int_{g\left( s\right) }^{g\left( x\right) }\left( g\left( t\right) -g\left(
s\right) \right) ^{N-\beta -1}dg\left( t\right) =  \tag{55}  \label{55}
\end{equation}
\begin{equation*}
\frac{\left \Vert \left( f\circ g^{-1}\right) ^{\left( N\right) }\circ
g\right \Vert _{\infty ,\mathbb{R}}}{\Gamma \left( N-\beta +1\right) }\left(
g\left( x\right) -g\left( s\right) \right) ^{N-\beta }.
\end{equation*}
That is
\begin{equation}
\left \vert D_{x-;g}^{\beta }f\left( s\right) \right \vert \leq \frac{\left
\Vert \left( f\circ g^{-1}\right) ^{\left( N\right) }\circ g\right \Vert
_{\infty ,\mathbb{R}}}{\Gamma \left( N-\beta +1\right) }\left( g\left(
x\right) -g\left( s\right) \right) ^{N-\beta },  \tag{56}  \label{56}
\end{equation}
$\forall $ $s\in (-\infty ,x].$

Clearly then
\begin{equation}
D_{x-;g}^{\beta }f\left( x\right) =0.  \tag{57}  \label{57}
\end{equation}
\end{proof}

As an application of Theorem \ref{t11} we get:

\begin{theorem}
\label{t12} All as in Theorem \ref{t11}, and under the assumption that $%
D_{\ast x;g}^{\beta }f\left( t\right) ,$ $D_{x-;g}^{\beta }f\left( t\right) $
are both continuous and bounded in $\left( x,t\right) \in \mathbb{R}^{2}$; $%
n\in \mathbb{N}:n\geq 2$. Then

1) if $N=1$ we have
\begin{equation*}
\underset{i\neq j\text{ or }i=j}{\underset{i,j\in \left \{ 1,...,7\right \} ,%
}{\max }}\left \Vert A_{n}^{\left( i,j\right) }\left( f\right) -f\right
\Vert _{\infty ,\left[ -1,1\right] }\leq \frac{\omega _{1}^{\beta }\left( g,%
\frac{1}{n^{1-\alpha }}\right) }{\Gamma \left( \beta +1\right) }
\end{equation*}
\begin{equation}
\left[ \underset{x\in \left[ -1,1\right] }{\sup }\omega _{1}\left(
D_{x-;g}^{\beta }f;\frac{1}{n^{1-\alpha }}\right) _{(-\infty ,x]}+\underset{%
x\in \left[ -1,1\right] }{\sup }\omega _{1}\left( D_{\ast x;g}^{\beta }f;%
\frac{1}{n^{1-\alpha }}\right) _{[x,+\infty )}\right] ,  \tag{58}  \label{58}
\end{equation}

2) in general we have
\begin{equation}
\underset{i\neq j\text{ or }i=j}{\underset{i,j\in \left \{ 1,...,7\right \} ,%
}{\max }}\left \Vert A_{n}^{\left( i,j\right) }\left( f\right) -f\right
\Vert _{\infty ,\left[ -1,1\right] }\leq \sum_{j_{\ast }=1}^{N-1}\frac{\left
\Vert \left( f\circ g^{-1}\right) ^{\left( j_{\ast }\right) }\circ g\right
\Vert _{\infty ,\left[ -1,1\right] }}{j_{\ast }!}\omega _{1}^{j_{\ast
}}\left( g,\frac{1}{n^{1-\alpha }}\right)  \tag{59}  \label{59}
\end{equation}
\begin{equation*}
+\frac{\omega _{1}^{\beta }\left( g,\frac{1}{n^{1-\alpha }}\right) }{\Gamma
\left( \beta +1\right) }\left[ \underset{x\in \left[ -1,1\right] }{\sup }%
\omega _{1}\left( D_{x-;g}^{\beta }f;\frac{1}{n^{1-\alpha }}\right)
_{(-\infty ,x]}+\underset{x\in \left[ -1,1\right] }{\sup }\omega _{1}\left(
D_{\ast x;g}^{\beta }f;\frac{1}{n^{1-\alpha }}\right) _{[x,+\infty )}\right]
.
\end{equation*}

An interesting case is when $\beta =\frac{1}{2}.$

Here we get generalized fractionally with rates the uniform convergence of $%
A_{n}^{\left( i,j\right) }\left( f\right) \rightarrow f$, over $\left[ -1,1%
\right] $, as $n\rightarrow +\infty $; $i,j\in \left \{ 1,...,7\right \} .$
\end{theorem}

\begin{proof}
By Theorem \ref{t11}, and by
\begin{equation*}
\left \Vert D_{\ast x;g}^{\beta }f\right \Vert _{\infty }\leq K_{1},\text{ \ }
\forall \text{ }x\in \mathbb{R};
\end{equation*}
and
\begin{equation*}
\left \Vert D_{x-;g}^{\beta }f\right \Vert _{\infty }\leq K_{2},\text{ \ }%
\forall \text{ }x\in \mathbb{R},
\end{equation*}
where $K_{1},K_{2}>0.$

We get that
\begin{equation*}
\omega _{1}\left( D_{x-;g}^{\beta }f;\xi \right) _{(-\infty ,x]}\leq 2K_{2},
\end{equation*}
and
\begin{equation*}
\omega _{1}\left( D_{\ast x;g}^{\beta }f;\xi \right) _{[x,+\infty )}\leq
2K_{1},\text{ \ }\forall \text{ }x\in \mathbb{R},
\end{equation*}
for any $\xi \geq 0.$
\end{proof}

It follows

\begin{corollary}
\label{c13} Let $\beta >0$, $\beta \notin \mathbb{N}$, $N=\left \lceil \beta
\right \rceil $. Let also a strictly increasing $g\in C^{1}\left( \mathbb{R}%
\right) $, such that $g^{-1}\in C^{N}\left( \mathbb{R}\right) $; $g$ is also
assumed to be uniformly continuous. Let also $f\in C^{N}\left( \mathbb{R}%
\right) $ and assume that $\left( f\circ g^{-1}\right) ^{\left( N\right)
}\circ g\in L_{\infty }\left( \mathbb{R}\right) $. Let $x\in \left[ -\phi
,\phi \right] $, $\phi >0$ and $n\in \mathbb{N}:n\geq 1+\phi ,$ $0<\alpha <1$
. Here $D_{\ast x;g}^{\beta }f\left( t\right) $, $D_{x-;g}^{\beta }f\left(
t\right) $ are both continuous and bounded in $\left( x,t\right) \in \mathbb{%
R}^{2}$. Then

1) If $N=1,$ we have
\begin{equation}
\underset{i\neq j\text{ or }i=j}{\underset{i,j\in \left \{ 1,...,7\right \} ,%
}{\max }}\left \vert A_{n}^{\left( i,j\right) }\left( f\right) \left(
x\right) -f\left( x\right) \right \vert \leq \frac{\omega _{1}^{\beta
}\left( g,\frac{1}{n^{1-\alpha }}\right) }{\Gamma \left( \beta +1\right) }
\tag{60}  \label{60}
\end{equation}
\begin{equation*}
\left[ \underset{x\in \left[ -\phi ,\phi \right] }{\sup }\omega _{1}\left(
D_{x-;g}^{\beta }f;\frac{1}{n^{1-\alpha }}\right) _{(-\infty ,x]}+\underset{%
x\in \left[ -\phi ,\phi \right] }{\sup }\omega _{1}\left( D_{\ast
x;g}^{\beta }f;\frac{1}{n^{1-\alpha }}\right) _{[x,+\infty )}\right] ,
\end{equation*}

2) in general we have
\begin{equation*}
\underset{i\neq j\text{ or }i=j}{\underset{i,j\in \left \{ 1,...,7\right \} ,%
}{\max }}\left \vert A_{n}^{\left( i,j\right) }\left( f\right) \left(
x\right) -f\left( x\right) \right \vert \leq \sum_{j_{\ast }=1}^{N-1}\frac{%
\left \vert \left( f\circ g^{-1}\right) ^{\left( j_{\ast }\right) }\left(
g\left( x\right) \right) \right \vert }{j_{\ast }!}\omega _{1}^{j_{\ast
}}\left( g,\frac{1}{n^{1-\alpha }}\right)
\end{equation*}
\begin{equation}
+\frac{\omega _{1}^{\beta }\left( g,\frac{1}{n^{1-\alpha }}\right) }{\Gamma
\left( \beta +1\right) }\left[ \underset{x\in \left[ -\phi ,\phi \right] }{%
\sup }\omega _{1}\left( D_{x-;g}^{\beta }f;\frac{1}{n^{1-\alpha }}\right)
_{(-\infty ,x]}+\underset{x\in \left[ -\phi ,\phi \right] }{\sup }\omega
_{1}\left( D_{\ast x;g}^{\beta }f;\frac{1}{n^{1-\alpha }}\right)
_{[x,+\infty )}\right] .  \tag{61}  \label{61}
\end{equation}
\end{corollary}

\begin{proof}
By Theorem \ref{t11}.
\end{proof}

By Corollary \ref{c13}, and monotonicity and subadditivity properties of $%
\left \Vert \cdot \right \Vert _{p}$, $p\geq 1$, we obtain the following $%
L_{p} $ result.

\begin{corollary}
\label{c14} Here all as in Corollary \ref{c13} and $p\geq 1$. Then

1) If $N=1,$ we have
\begin{equation}
\underset{i\neq j\text{ or }i=j}{\underset{i,j\in \left \{ 1,...,7\right \} ,%
}{\max }}\left \Vert A_{n}^{\left( i,j\right) }f-f\right \Vert _{p,\left[
-\phi ,\phi \right] }\leq 2^{\frac{1}{p}}\phi ^{\frac{1}{p}}\frac{\omega
_{1}^{\beta }\left( g,\frac{1}{n^{1-\alpha }}\right) }{\Gamma \left( \beta
+1\right) }  \tag{62}  \label{62}
\end{equation}
\begin{equation*}
\left[ \underset{x\in \left[ -\phi ,\phi \right] }{\sup }\omega _{1}\left(
D_{x-;g}^{\beta }f;\frac{1}{n^{1-\alpha }}\right) _{(-\infty ,x]}+\underset{%
x\in \left[ -\phi ,\phi \right] }{\sup }\omega _{1}\left( D_{\ast
x;g}^{\beta }f;\frac{1}{n^{1-\alpha }}\right) _{[x,+\infty )}\right] ,
\end{equation*}
and

2) in general we get
\begin{equation}
\underset{i\neq j\text{ or }i=j}{\underset{i,j\in \left \{ 1,...,7\right \} ,%
}{\max }}\left \Vert A_{n}^{\left( i,j\right) }f-f\right \Vert _{p,\left[
-\phi ,\phi \right] }\leq \sum_{j_{\ast }=1}^{N-1}\frac{\left \Vert \left(
f\circ g^{-1}\right) ^{\left( j_{\ast }\right) }\circ g\right \Vert _{p,%
\left[ -\phi ,\phi \right] }}{j_{\ast }!}\omega _{1}^{j_{\ast }}\left( g,%
\frac{1}{n^{1-\alpha }}\right)  \tag{63}  \label{63}
\end{equation}
\begin{equation*}
+2^{\frac{1}{p}}\phi ^{\frac{1}{p}}\frac{\omega _{1}^{\beta }\left( g,\frac{1%
}{n^{1-\alpha }}\right) }{\Gamma \left( \beta +1\right) }
\end{equation*}
\begin{equation*}
\left[ \underset{x\in \left[ -\phi ,\phi \right] }{\sup }\omega _{1}\left(
D_{x-;g}^{\beta }f;\frac{1}{n^{1-\alpha }}\right) _{(-\infty ,x]}+\underset{%
x\in \left[ -\phi ,\phi \right] }{\sup }\omega _{1}\left( D_{\ast
x;g}^{\beta }f;\frac{1}{n^{1-\alpha }}\right) _{[x,+\infty )}\right] .
\end{equation*}
By (\ref{62}), (\ref{63}) we derive the generalized fractional $L_{p}$, $%
p\geq 1$, convergence with rates of $A^{\left( i,j\right) }f$ to $f$, all $%
i,j\in \left \{ 1,...,7\right \} .$
\end{corollary}

\begin{thebibliography}{99}
\bibitem{1} C.D. Aliprantis, K.C. Border, \textit{Infinite Dimensional
Analysis},\ Springer, New York, 2006.

\bibitem{2} G.A. Anastassiou, \textit{Rate of Convergence of some neural
network operators to the unit - univariate case}, Journal of Mathematical
Analysis and Applications, Vol 212 (1997), 237-262.

\bibitem{3} G.A. Anastassiou, \textit{On right fractional calculus}, Chaos,
Solitons and Fractals, Vol. 42 (209), 365-376.

\bibitem{4} G.A. Anastassiou, \textit{Intelligent Systems II: Complete
Approximation by Neural Network Operators}, Springer, Heidelberg, New York,
2016.

\bibitem{5} G.A. Anastassiou, \textit{Principles of General Fractional
Analysis for Banach space valued functions}, Bull. Allahabad Math. Soc.,
32(1) (2017), 71-145.

\bibitem{6} G.A. Anastassiou, \textit{Banach Space Valued Neural Network},
Springer, Heidelberg, New York, 2023.

\bibitem{7} G.A. Anastassiou, \textit{Parametrized, Deformed and General
Neural Networks}, Springer, Heidelberg, New York, 2023.

\bibitem{8} P. Cardaliaguet and G. Euvrard, \textit{Approximation of a
function and its derivative with a neural network}, Neural Networks, Vol. 5
(1992), 207-220.

\bibitem{9} K. Diethelm, \textit{The Analysis of Fractional Differential
Equations}, Lecture Notes in Mathematics 2004, Springer-Verlag, Berlin,
Heidelberg, 2010.

\bibitem{10} A.M.A. El-Sayed and M. Gaber, \textit{On the finite Caputo and
finite Riesz derivatives}, Electronic Journal of Theoretical Physics, Vol.
3, No. 12 (2006), 81-95.

\bibitem{11} G.A. Anastassiou, \textit{Intelligent Computations: Abstract
Fractional Calculus, Inequalities, Approximations}, Springer, Heidelberg,
New York, 2018.

\bibitem{12} Z. Chen and F. Cao, \textit{The approximation operators with
sigmoidal functions}, Computers and Mathematics with Applications, 58
(2009), 758-765.

\bibitem{13} D. Costarelli, R. Spigler, \textit{Approximation results for
neural network operators activated by sigmoidal functions}, Neural Networks
44 (2013), 101-106.

\bibitem{14} D. Costarelli, R. Spigler, \textit{Multivariate neural network
operators with sigmoidal activation functions}, Neural Networks 48 (2013),
72-77.

\bibitem{15} S. Haykin, \textit{Neural Networks: A Comprehensive Foundation}
(2 ed.), Prentice Hall, New York, 1998.

\bibitem{16} W. McCulloch and W. Pitts, \textit{A logical calculus of the
ideas immanent in nervous activity}, Bulletin of Mathematical Biophysics, 7
(1943), 115-133.

\bibitem{17} J. Mikusinski, \textit{The Bochner Integral}, Academic Press,
New York, 1978.

\bibitem{18} T.M. Mitchell, \textit{Machine Learning}, WCB-McGraw-Hill, New
York, 1997.

\bibitem{19} D. S. Yu; F. L. Cao, \textit{Construction and approximation
rate for feed-forward neural network operators with sigmoidal functions}, J.
Comput. Appl. Math. 453 (2025), paper No. 116150.

\bibitem{20} Cen, Siyu; Jin, Bangti; Quan, Qimeng; Zhou, Zhi; \textit{Hybrid
neural-network FEM approximation of diffusion coefficient in elliptic and
parabolic problems}. IMA J. Numer. Anal. 44 (2024), no. 5, 3059-3093.

\bibitem{21} Coroianu Lucian; Costarelli, Danillo; Natale, Mariarosaria;
Panti\c{s}, Alexandra; \textit{The approximation capabilities of
Durrmeyer-type neural network operators}. J. Appl. Math. Comput. 70 (2024),
no. 5, 4581-4599.

\bibitem{22} Warin, Xavier; \textit{The GroupMax neural network
approximation of convex functions}. IEEE Trans. Neural Netw. Learn. Syst. 35
(2024), no. 8, 11608-11612.

\bibitem{23} Fabra, Arnau; Guasch, Oriol; Baiges, Joan; Codina, Ramon;
\textit{Approximation of acoustic black holes with finite element mixed
formulations and artificial neural network correction terms}. Finite Elem.
Anal. Des. 241 (2024), paper No. 104236.

\bibitem{24} Grohs, Philipp; Voigtlaender, Felix. \textit{Proof of the
theory-to-practice gap in deep learning via sampling complexity bounds for
neural network approximation spaces}. Found. Comput. Math. 24 (2024), no. 4,
1085-1143.

\bibitem{25} Basteri, Andrea; Trevisan, Dario; \textit{Quantitative Gaussian
approximation of randomly initialized deep neural networks}. Mach. Learn.
113 (2024), no. 9, 6373-6393.

\bibitem{26} De Ryck, T.; Mishra, S. \textit{Error analysis for deep neural
network approximations of parametric hyperbolic conservation laws}. Math.
Comp. 93 (2024), no. 350, 2643-2677.

\bibitem{27} Liu, Jie; Zhang, Baoji; Lai, Yuyang; Fang, Liqiau. \textit{Hull
form optimization research based on multi-precision back-propagation neural
network approximation model}. Internal. J. Numer. Methods Fluid 96 (2024),
no. 8, 1445-1460.

\bibitem{28} Yoo, Jihahm; Kim, Jaywon; Gim, Minjung; Lee, Haesung. \textit{%
Error estimates of physics-informed neural networks for initial value
problems}. J. Korean Soc. Ind. Appl. Math. 28 (2024), no. 1, 33-58.
\end{thebibliography}

\end{document}
